\chapter{Felhantering}
\label{ch:fel}

{\large\itshape Fem sorters fel, en signal, och en nod som stänger av sig själv.\par}

\begin{chapterbox}{I det här kapitlet}
\begin{itemize}
  \item De fem fel en CAN-nod kan upptäcka, och vem som upptäcker vad.
  \item Felframen: hur en nod avbryter en pågående överföring.
  \item Felräknarna, och de tre tillstånden en nod kan vara i.
  \item Bus-off: när en nod kopplar bort sig själv, och varför det är rätt beteende.
\end{itemize}
\end{chapterbox}

\section{Fem sorters fel}

CAN upptäcker fel på fem sätt. Tre av dem har vi redan gått igenom under andra namn:

\begin{booktable}{@{}lLl@{}}
\thead{Fel} & \thead{Vad det är} & \thead{Vem ser det} \\ \midrule
Bitfel & En sändande nod läser tillbaka något annat än den sände & Sändaren \\
Stoppfel & Sex lika bitar i rad där femregeln kräver en flank & Alla \\
CRC-fel & Den mottagna checksumman stämmer inte & Mottagare \\
Formatfel & Ett fält som måste vara recessivt är det inte & Alla \\
Kvittensfel & Ingen drog ned ACK-luckan & Sändaren \\
\end{booktable}

Tillsammans ger de en restfelssannolikhet som brukar anges till omkring $10^{-11}$ av de fel som
uppstår --- alltså ungefär ett oupptäckt fel per tusen år på en hårt belastad buss. Det är inte
en slump utan resultatet av att fem oberoende mekanismer överlappar.

\subsection{Bitfel, och undantagen}

En sändande nod jämför hela tiden det den sänder med det som ligger på bussen. Skiljer de sig är
det ett bitfel --- \emph{utom} på två ställen där en skillnad är fullt normal:

\begin{itemize}
  \item \textbf{Under arbitreringsfältet.} Recessiv ut och dominant in betyder förlorad
    arbitrering, inte fel (\kapref{ch:arbitrering}).
  \item \textbf{I ACK-luckan.} Sändaren sänder recessivt och \emph{hoppas} på dominant
    (\kapref{ch:frameformat}).
\end{itemize}

Att skilja de två fallen från ett riktigt bitfel är en av de saker som gör en CAN-kontroller
mer invecklad än den ser ut.

\section{Felframen}

När en nod upptäcker ett fel säger den inte till efteråt. Den avbryter framen \textbf{medan den
pågår}, genom att sända en \textbf{felframe}: sex dominanta bitar i rad.

Sex dominanta bitar är omöjliga i en korrekt frame --- femregeln ser till att det aldrig går mer
än fem --- så mönstret kan inte förväxlas med data. Varje annan nod på bussen ser det, förstår att
framen är trasig, och kastar den.

Och eftersom alla noder ser felet kommer flera av dem att sända sin egen felflagga i sin tur.
Resultatet är en serie dominanta bitar, följd av en avgränsare på åtta recessiva.

Två saker följer:

\textbf{Alla noder kastar framen, inte bara den som upptäckte felet.} Bussen hålls konsistent:
antingen tar alla emot framen, eller ingen. Det är en starkare garanti än den ett
hemmasnickrat protokoll normalt ger sig på.

\textbf{Sändaren skickar om automatiskt.} Kontrollern gör det själv, utan att mjukvaran blir
inblandad. En drivrutin som skickar en frame och får veta att det gick bra kan alltså ha fått
den omsänd flera gånger utan att märka det.

\section{Felräknarna}

Här kommer den del av CAN som är minst uppenbar och mest genomtänkt.

En nod som råkar ut för fel ska säga ifrån. En nod som är \emph{trasig} ska sluta säga ifrån ---
annars förstör den bussen för alla andra genom att sända felframes i ett.

CAN skiljer dem åt med två räknare i varje nod:

\begin{itemize}
  \item \textbf{TEC} --- transmit error counter, fel medan noden sände.
  \item \textbf{REC} --- receive error counter, fel medan den tog emot.
\end{itemize}

Reglerna i sammandrag: ett fel ökar räknaren med 8 (sändarsidan) eller 1 (mottagarsidan), och en
lyckad överföring minskar den med 1.

\textbf{Asymmetrin är hela poängen.} En nod som verkligen är trasig är inblandad i varje fel och
ökar snabbt. En nod som bara råkar vara med på en störig buss ökar långsamt och hinner arbeta av
sig. Räknarna mäter alltså inte "hur många fel har hänt" utan "hur sannolikt är det att felet är
mitt".

\subsection{De tre tillstånden}

\bookfigure{error_states}{En nods tre feltillstånd och övergångarna mellan dem.}{fig:errorstates}

\begin{booktable}{@{}llL@{}}
\thead{Tillstånd} & \thead{Villkor} & \thead{Beteende} \\ \midrule
Error active & TEC och REC $\leq$ 127 & Normal drift; sänder \emph{aktiva} felflaggor (dominanta) \\
Error passive & Någon räknare $>$ 127 & Sänder \emph{passiva} felflaggor (recessiva) och väntar extra mellan sändningar \\
Bus off & TEC $>$ 255 & Sänder ingenting alls \\
\end{booktable}

Övergången till \emph{error passive} är den intressanta. En passiv felflagga är recessiv, vilket
betyder att den inte stör någon annan. Noden får fortfarande delta, den får fortfarande
rapportera, men den kan inte längre dra ned bussen för allas räkning.

\subsection{Bus off}

Vid TEC över 255 kopplar noden bort sig själv. Den sänder inte, kvitterar inte, stör inte.

Det är protokollets sista försvar, och det är värt att stanna vid hur ovanligt beteendet är: en
nod som konstaterar att den själv troligen är den trasiga delen, och som drar den slutsatsen
utan att någon talar om det för den.

En nod i bus off kan tas tillbaka --- typiskt efter att ha sett bussen ledig ett visst antal gånger,
eller efter att mjukvaran återinitierar kontrollern. Att den återhämtningen finns är också
avsiktligt: felet kan ha varit en lös kontakt som nu sitter.

\begin{aside}
\note[För den som skriver drivrutin] Bus off är det ena tillstånd mjukvaran inte kan ignorera.
Allt annat hanterar kontrollern själv --- omsändningar, felframes, räknare --- men en nod i bus off
kommer aldrig tillbaka av sig själv om konstruktionen kräver en återinitiering. Det är ett av de
få ställen där felhanteringen når ända upp i applikationen.
\end{aside}

\section{Vad mjukvaran ser}

Lägg märke till hur lite av allt det här som når mjukvaran i en fullständig kontroller:

\begin{itemize}
  \item Fel upptäcks i hårdvara.
  \item Felframes sänds i hårdvara.
  \item Omsändning sker i hårdvara.
  \item Räknarna uppdateras i hårdvara.
\end{itemize}

Mjukvaran får en statusbit, kanske en felkod, och möjligheten att läsa räknarna. Det är en
enorm mängd arbete som gjorts åt den --- jämför med ACK, timeout, retry och dubbletthantering i ett
hemsnickrat protokoll, som är samma problem löst i mjukvara, sämre.

\begin{simplification}
\note[Förenkling] Den förenklade kontrollern i \kapref{ch:kursen} implementerar \emph{ingenting}
av det här kapitlet utom stopp\-felet och CRC-kontrollen. Den sänder inga felframes, har inga
felräknare, känner inte till error passive eller bus off, och skickar aldrig om automatiskt.
Den signalerar bara att något gick fel, med en enda bit, och lämnar resten till mjukvaran.
\end{simplification}

\section*{Övningar}

\exercise{Vem ser felet?}{Resonemang}
\label{ex:6:1}
För vart och ett av följande: vilken sorts fel är det, och vem upptäcker det?
\begin{enumerate}[a)]
  \item En störning flippar en bit i datafältet.
  \item En störning flippar en bit i ett fält som måste vara recessivt.
  \item Sändaren är ensam på bussen.
  \item Sju lika bitar i rad passerar i arbitreringsfältet.
\end{enumerate}

\exercise{Räknarna}{Räkning}
\label{ex:6:2}
En nod börjar på TEC 0 och råkar ut för åtta sändningsfel i rad, följt av tjugo lyckade
sändningar.
\begin{enumerate}[a)]
  \item Vad står TEC på efter felen?
  \item Efter de lyckade sändningarna?
  \item I vilket tillstånd är noden vid respektive tidpunkt?
\end{enumerate}

\exercise{Asymmetrin}{Resonemang}
\label{ex:6:3}
Ett fel ökar sändarens räknare med 8 men mottagarens med 1.
\begin{enumerate}[a)]
  \item Varför är ökningen olika stor?
  \item En nod med trasig sändare och en nod på en störig buss ser likadana ut i början. Vad
    skiljer dem åt efter en stund?
\end{enumerate}

\exercise{Bus off}{Resonemang}
\label{ex:6:4}
\begin{enumerate}[a)]
  \item Varför kopplar en nod bort sig själv i stället för att fortsätta rapportera fel?
  \item Vad hade hänt med bussen om den inte gjorde det?
  \item Varför finns en väg tillbaka från bus off?
\end{enumerate}
